\documentclass{article}

\usepackage{amsmath,amssymb}
\usepackage{xurl}
\usepackage[hidelinks]{hyperref}
\setlength{\emergencystretch}{2em}

% Shared commands for notes in docs/paper-gaps/.

\DeclareUnicodeCharacter{2080}{\ifmmode{}_{0}\else\textsubscript{0}\fi}
\DeclareUnicodeCharacter{2081}{\ifmmode{}_{1}\else\textsubscript{1}\fi}
\DeclareUnicodeCharacter{2082}{\ifmmode{}_{2}\else\textsubscript{2}\fi}
\DeclareUnicodeCharacter{2099}{\ifmmode{}_{n}\else\textsubscript{n}\fi}

\newcommand{\Lean}{\textsc{Lean}}
\ExplSyntaxOn
\NewDocumentCommand{\leanid}{m}{
  \ifmmode
    \textsf{\detokenize{#1}}
  \else
    \group_begin:
    \urlstyle{sf}
    \tl_set:Nn \l_tmpa_tl {#1}
    \tl_replace_all:Nnn \l_tmpa_tl {\_} {_}
    \exp_args:NV \path \l_tmpa_tl
    \group_end:
  \fi
}
\ExplSyntaxOff
\providecommand{\tr}{\operatorname{tr}}
\newcommand{\tnleanIssueBase}{https://github.com/LionSR/TNLean/issues/}
\newcommand{\tnleanPrBase}{https://github.com/LionSR/TNLean/pull/}
\newcommand{\tnleanDocBase}{https://sirui-lu.com/TNLean/docs/}
\newcommand{\ghissue}[1]{\href{\tnleanIssueBase#1}{GitHub issue \##1}}
\newcommand{\ghpr}[1]{\href{\tnleanPrBase#1}{PR \##1}}
\newcommand{\doclink}[1]{\href{\tnleanDocBase#1.html}{\path{#1}}}

% Dirac notation and the identity operator, as in blueprint/src/macros/common.tex.
% \providecommand keeps note-local definitions authoritative.
\usepackage{dsfont}
\providecommand{\ket}[1]{|#1\rangle}
\providecommand{\bra}[1]{\langle #1|}
\providecommand{\braket}[2]{\langle #1 | #2 \rangle}
\providecommand{\ketbra}[2]{|#1\rangle\!\langle #2|}
\providecommand{\Id}{\mathds{1}}
\providecommand{\one}{\mathds{1}}
\providecommand{\C}{\mathbb{C}}
\providecommand{\MN}[1]{M_{#1}(\C)}
\providecommand{\MD}{M_D(\C)}

% External referencing for paper sources.  The build helper writes sanitized
% auxiliary files under build/paper-aux/ and excludes duplicated source labels.
\usepackage{xr}

% Import a generated paper auxiliary file with a caller-supplied label prefix.
% The candidate paths support builds from docs/paper-gaps/, the repository root,
% and build/paper-aux/ itself.
\newcommand{\importpaperaux}[2]{%
  \IfFileExists{../../build/paper-aux/#2.aux}{%
    \externaldocument[#1]{../../build/paper-aux/#2}%
  }{%
    \IfFileExists{build/paper-aux/#2.aux}{%
      \externaldocument[#1]{build/paper-aux/#2}%
    }{%
      \IfFileExists{#2.aux}{%
        \externaldocument[#1]{#2}%
      }{}%
    }%
  }%
}

% General external reference macros
\newcommand{\paperref}[2]{\ref{#1-#2}}
\newcommand{\papereqref}[2]{(\ref{#1-#2})}

% CPSV16 source (1606.00608 / MPDO-22-12-17-2.tex) external document and shortcuts
\importpaperaux{cpsv16-}{1606.00608/MPDO-22-12-17-2-tnlean}
\newcommand{\cpsvref}[1]{\paperref{cpsv16}{#1}}
\newcommand{\cpsveqref}[1]{\papereqref{cpsv16}{#1}}

% Verdict marker: \gapnote{<kind>}{<status>}. Kind is the policy
% classification (clarification, local-correction, scope-restriction,
% unfaithful, false-source, open-gap); status is open, wip (elimination
% underway), resolved, or historical. Machine-read by the site index;
% prints a verdict line.
\newcommand{\gapnote}[2]{\par\noindent\textbf{Verdict:} #1 (#2).\par}


\title{Zero Multiplicities in Wolf's Two-Pure-State Spectrum Shorthand}
\author{}
\date{2026-08-08}

\begin{document}
\maketitle

\section{The source display}

In the proof of the spectrum-preserver consequence following Wigner's theorem,
Wolf considers Hermitian rank-one projections $P,Q\in M_d(\mathbb C)$ and
writes that the spectrum of $P+Q$ is
\begin{align}
  1\pm\sqrt{\tr(PQ)}.
\end{align}
The statement appears in Chapter~1, lines 289--296 of the local source,
\path{Notes/WolfNoteTexSource/ch01_deconstructing_quantum.tex}.

\section{The suppressed eigenvalues}

The range of $P+Q$ lies in the span of the two one-dimensional ranges of $P$
and $Q$. Hence $P+Q$ vanishes on a subspace of dimension at least $d-2$.
Counting algebraic multiplicities, its characteristic polynomial is
\begin{align}
  \chi_{P+Q}(X)
  =X^{d-2}((X-1)^2-\tr(PQ))
  \qquad(d\geq2).
\end{align}
Thus the displayed pair gives the two eigenvalues arising on the span of the
pure states, but it suppresses the remaining $d-2$ zero eigenvalues and their
multiplicities. In dimension one the separate formula is
\begin{align}
  \chi_{P+Q}(X)=X-2,
\end{align}
because every normalized rank-one projection is the identity. In dimension
zero there are no projective pure states.

\section{Formal correction}

The declaration
\leanid{Projectivization.charpoly_add_pureStateMatrix_of_two_le} proves the
multiplicity-aware formula through a $d\times2$ and $2\times d$ factorization.
The declaration
\leanid{Projectivization.trace_pureStateMatrix_mul_eq_of_charpoly_add_eq}
then recovers $\tr(PQ)$ from equality of two such characteristic polynomials.
Both are in
\doclink{TNLean/Channel/Wigner/TwoPureStateCharpoly}; the low-dimensional cases
are stated explicitly there. This correction is tracked by \ghissue{5683}.

\section{Classification}

This is a local correction to a spectrum shorthand, not a claim that Wolf's
intended spectrum-preserver argument is false. The displayed pair identifies
the nonzero two-dimensional part used to recover the transition probability.
The formal statement uses characteristic polynomials because they retain all
zero eigenvalues and all algebraic multiplicities without ambiguity.

\bibliographystyle{alpha}
\bibliography{references}

\end{document}